% !TEX program = xelatex
\documentclass[10pt, letterpaper]{article}

% Packages:
\usepackage[
    ignoreheadfoot, % set margins from page edge to text
    top=1.5 cm,
    bottom=1.5 cm,
    left=1.5 cm,
    right=1.5 cm,
    footskip=1.0 cm,
    % showframe, % search for "showframe" for debugging
]{geometry}
\usepackage{titlesec}
\usepackage{tabu}
\usepackage{enumitem}
\usepackage[hidelinks]{hyperref}
\usepackage{fancyhdr}
\usepackage[english]{babel}
\usepackage{bookmark}
\usepackage{xcolor}
\usepackage{fontspec}

% Font:
\defaultfontfeatures{Ligatures=TeX}
\IfFontExistsTF{Times New Roman}
  {\setmainfont{Times New Roman}}
  {\setmainfont{TeX Gyre Termes}}
\linespread{1.03}\selectfont

% Formatting:
\pagestyle{fancy}
\fancyhf{} 
\renewcommand{\headrulewidth}{0pt}
\renewcommand{\footrulewidth}{0pt}
\urlstyle{same}
\setlength{\parindent}{0pt}

% Custom commands
\newcommand{\entry}[4]{
    \textbf{#1} \hfill #2 \\
    \textit{#3} \hfill \textit{#4}
}

\newcommand{\myname}{\textbf{\textup{Junhoo Lee}}}
\newcommand{\coauthors}[1]{\textit{#1}}
\newcommand{\pubid}[1]{\textbf{[#1]}}
\newenvironment{publicationlist}{
    \begin{description}[leftmargin=4.0em, labelwidth=3.2em, labelsep=0.6em, itemsep=0.85mm, topsep=0.5mm, style=multiline, font=\normalfont]
}{
    \end{description}
}
\newcommand{\oralpresentation}[1]{{\color{red}\textbf{Oral Presentation (#1)}}}

\newcommand{\selpub}[5]{
    \item \textbf{#1}, \textbf{#2} {\footnotesize\textit{\textcolor{gray}{#3}}} \\[0.2mm]
    #4 \\[0.15mm]
    {\footnotesize #5}
}

\newcommand{\projectentry}[4]{
    \textbf{#1} \hfill {\scriptsize #2} \\
    {\scriptsize\textit{#3}} \\
    {\scriptsize $\bullet$ #4} \\
    \vspace{0.9mm}
}

\newcommand{\talkentry}[4]{
    \textbf{\href{#1}{\textcolor{blue}{#2}}} \hfill {\scriptsize #3} \\
    {\scriptsize\textit{#4}}
}

\newcommand{\skillentry}[2]{
    \textbf{#1}: #2\par
}

\titleformat{\section}
  {\large\bfseries\scshape\raggedright\uppercase}
  {}{0em}
  {}
  [\titlerule]

\titlespacing{\section}{0pt}{6pt}{3pt}

\begin{document}

% Header
\begin{center}
    {\Huge \textbf{JUNHOO LEE}} \\ \vspace{2mm}
    Ph.D. Candidate, Seoul National University \\ \vspace{1mm}
    \small
    \href{mailto:mrjunoo@snu.ac.kr}{mrjunoo@snu.ac.kr} $|$ \href{https://junhoo.me}{https://junhoo.me}
\end{center}

% Education
\section{Education}
\textbf{Seoul National University} \hfill Seoul, South Korea \\
\textit{Ph.D. Candidate in Intelligence and Information} \hfill Sep 2021 -- Present (Exp. Aug 2026)
\newline
{\small Advisor: Prof. Nojun Kwak, Machine Intelligence \& Pattern Analysis Laboratory (MIPAL)}

\vspace{1mm}

\textbf{Seoul National University} \hfill Seoul, South Korea \\
\textit{B.Sc. in Electrical and Computer Engineering} \hfill Mar 2017 -- Sep 2021
\newline
{\small B.S. Thesis: \textit{Improvement of Elevator Control Algorithm Using Reinforcement Learning}. Advisor: Sungroh Yoon}

% Research Interests
\section{Research Interests}
{\normalsize
My research studies how foundation models acquire reusable structure, and how
that structure can be controlled, adapted, and diagnosed across language, vision,
generative modeling, and embodied interfaces. My work on discrete diffusion
language models introduces inference-time planning structure for controllable
generation [C1], while my meta-learning papers study fast adaptation and task
generalization through optimization-trajectory structure and flexible episodic
interfaces [C6, C7]. I also work on model diagnosis, semantic robustness, and
data-efficient transfer, including deep support vectors for pretrained-model
interpretation [C5], black-box semantic fingerprinting for generative models
[C2], multimodal temporal grounding [C4], and instruction-following failures in
text-to-image generation [W1]. Recently, I have been extending this line toward
vision-language-action and embodied foundation models, where pretrained semantic
structure must remain reliable under changing instructions, visual observations,
and action interfaces.
}

% Publications
\section{Publications}

{\small
\begin{publicationlist}
    \item[\pubid{C1}] \textbf{\href{https://arxiv.org/abs/2510.13870}{Unlocking the Potential of Diffusion Language Models through Template Infilling}} \\
    \textit{\myname, Seungyeon Kim, Nojun Kwak} \\
    \textit{Annual Meeting of the Association for Computational Linguistics (\textbf{ACL})}, 2026. \\
    \oralpresentation{485/10518=4.61\%}
    \item[\pubid{C2}] \textbf{\href{https://junhoo.me/csf/csf-paper.pdf}{CSF: Black-box Fingerprinting via Compositional Semantics for Text-to-Image Models}} \\
    \textit{\myname, Mijin Koo, Nojun Kwak} \\
    \textit{The IEEE/CVF Conference on Computer Vision and Pattern Recognition (\textbf{CVPR})}, 2026.
    \item[\pubid{C3}] \textbf{\href{https://arxiv.org/abs/2510.13865}{Deep Edge Filter}} (Co-first) \\
    \textit{Dongkwan Lee$^\dagger$, \textbf{Junhoo Lee}$^\dagger$, Nojun Kwak} \\
    \textit{The Conference on Neural Information Processing Systems (\textbf{NeurIPS})}, 2025.
    \item[\pubid{C4}] \textbf{\href{https://arxiv.org/abs/2507.04667}{What's Making That Sound Right Now? Video-centric Audio-Visual Localization}} \\
    \textit{Hahyeon Choi, \myname, Nojun Kwak} \\
    \textit{The IEEE/CVF International Conference on Computer Vision (\textbf{ICCV})}, 2025.
    \item[\pubid{C5}] \textbf{\href{https://arxiv.org/abs/2403.17329}{Deep Support Vectors}} \\
    \textit{\myname, Hyunho Lee, Kyomin Hwang, Nojun Kwak} \\
    \textit{The Conference on Neural Information Processing Systems (\textbf{NeurIPS})}, 2024.
    \item[\pubid{C6}] \textbf{\href{https://arxiv.org/abs/2401.05097}{Any-Way Meta Learning}} \\
    \textit{\myname, Yearim Kim, Hyunho Lee, Nojun Kwak} \\
    \textit{The AAAI Conference on Artificial Intelligence (\textbf{AAAI})}, 2024.
    \item[\pubid{C7}] \textbf{\href{https://arxiv.org/abs/2310.02751}{SHOT: Suppressing the Hessian along the Optimization Trajectory}} \\
    \textit{\myname, Jayeon Yoo, Nojun Kwak} \\
    \textit{The Conference on Neural Information Processing Systems (\textbf{NeurIPS})}, 2023.
    \item[\pubid{W1}] \textbf{\href{https://arxiv.org/abs/2404.15154}{Do Not Think About Pink Elephant!}} (Co-first) \\
    \textit{Kyomin Hwang$^\dagger$, Suyoung Kim$^\dagger$, \textbf{Junhoo Lee}$^\dagger$, Nojun Kwak} \\
    \textit{CVPR Workshops -- Responsible Generative AI}, 2024.
    \item[\pubid{W2}] \textbf{Coreset Selection for Object Detection} \\
    \textit{Hojun Lee, Suyoung Kim, \myname, Jaeyoung Yoo, Nojun Kwak} \\
    \textit{CVPR Workshops -- Dataset Distillation}, 2024.
    \item[\pubid{W3}] \textbf{\href{https://arxiv.org/abs/2405.00348}{Practical Dataset Distillation Based on Deep Support Vectors}} \\
    \textit{Hyunho Lee, \myname, Nojun Kwak} \\
    \textit{ECCV Workshops -- Dataset Distillation Challenge}, 2024.
    \item[\pubid{J1}] \textbf{End-to-End Multi-Entity Customization} \\
    \textit{Wonhark Park$^*$, Jaehyun Lee$^*$, Wonsik Shin, \myname, Nojun Kwak} \\
    \textit{IET Image Processing}, Accepted. 2026.
    \item[\pubid{J2}] \textbf{\href{https://arxiv.org/abs/2508.20224}{The Role of Teacher Calibration in Knowledge Distillation}} \\
    \textit{Suyoung Kim, Seonguk Park, \myname, Nojun Kwak} \\
    \textit{IEEE Access}, 2025.
\end{publicationlist}
}

% Awards
\section{Awards \& Honors}
{\small
\begin{tabu} to 0.97\textwidth { X[0.8,l] X[0.2,r] }
    ICML Gold Reviewer & 2026 \\
    BK21 Future Innovation Talent Bronze Prize (USD 1,000) & 2023 \\
    BK21 Outstanding Research Talent Fellowship (USD 3,500) & 2023 \\
    Yulchon AI Star Scholarship (USD 8,000) & 2022 \\
    3rd Place, SNU FastMRI Challenge (USD 3,000) & 2021 \\
    Kwanak Scholarship (full funded) & 2021 \\
    National Science and Engineering Scholarship (full funded) & 2017--2018
\end{tabu}
}

% Research Projects
\section{Research Projects}
{\small
\textbf{Gradient-Descent-Based Meta-Learning Algorithms} (Samsung Electronics, Main Researcher, 2022--2023)\\
\textbf{Self-Directed Visual Object Detection for Shopping Assistance} (IITP, Project Manager, 2022--2024)\\
\textbf{AI Video Generation for Korean Traditional Media Art} (MCST/EASYWITH, Main Engineer, 2024--2026)\\
\textbf{Physically Plausible Video Diffusion Modeling} (AI Star Fellowship/MSIT, Project Manager, 2025).
}

% Skills
\section{Skills}
{\small
\textbf{Programming Languages}: Python, CUDA, R, C++, Verilog; \textbf{ML Development Stack}: PyTorch, TensorFlow, Slurm, Docker.
}

% Talks
\section{Talks}
{\small
\textbf{\href{https://junhoo.me/csf}{Compositional Semantic Fingerprinting (CSF) for Black-box Attribution of Text-to-Image Models}}, SWCS 2026, Apr. 2026\\
\textbf{\href{https://arxiv.org/abs/2403.17329}{Deep Support Vectors: Interpreting Deep Models via Support Vector Extraction}}, KCC 2024, Jun. 26, 2024.
}

% Academic Service
\section{Academic Service}
{\small
\textbf{Gold Reviewer}: ICML 2026.\\
\textbf{Conference Reviewer}: CVPR (2025, 2026), ICCV (2025), ECCV (2026), ICLR (2025, 2026), ICML (2026), NeurIPS (2024, 2025, 2026).
}

% Patents
\section{Patents}
{\small
\textbf{Method and Apparatus for Edge Filtering to Improve the Robustness of Deep Learning Models}, Korean Patent 10-2025-0215458\\
\textbf{Method and Apparatus for Controlling Unintended Object Generation in Image Generation Models through Prompt Reconfiguration}, Korean Patent 10-2025-0063788.
}

\end{document}
